%% ============================================================================
%%  demo.tex -- Muestrario de la clase lafrox
%%
%%  No es un entregable: es el catalogo visual de la plantilla. Sirve para
%%  (a) revisar el diseno de un golpe y (b) copiar y pegar la sintaxis de cada
%%  elemento al redactar los subdocumentos.
%%
%%  Compilar:  latexmk demo.tex        (dos pasadas, TikZ overlay)
%%  Limpiar:   latexmk -c
%% ============================================================================

% Opciones disponibles: borrador | final | firma | sinfirma | carta | oficio
%                       apa | sinnotas
\documentclass[carta,firma,borrador]{lafrox}

% ---------------------------------------------------------------------------
%  Metadatos de la oferta. Editar solo aqui: portada, encabezados, pies y
%  metadatos del PDF se alimentan de estas claves.
% ---------------------------------------------------------------------------
%  Las claves objeto, sobre, domicilio, correo, telefono y nomenclatura
%  alimentan el bloque que el Art. 49.1 exige leer en la caratula. Si se
%  omiten, la clase usa sus valores por defecto (seccion 6 de lafrox.cls),
%  que son PROVISIONALES mientras no se fije la identidad del PROPONENTE.
\datoslafrox{
  empresa       = LafroX SpA,
  rut           = 77.418.902-K,
  domicilio     = {Av. Brasil 2241, Valparaíso},
  representante = Alex Aravena,
  cargo         = Jefe de Proyecto y Apoderado,
  correo        = contacto@lafrox.cl,
  telefono      = +56 32 250 4100,
  cliente       = Distribuidora Puelche S.A.,
  objeto        = Proyecto de Plataforma Digital de Misión Crítica,
  sobre         = Sobre N.\textdegree{} 2 --- Oferta Técnica,
  documento     = Propuesta Técnica,
  subtitulo     = Muestrario de plantilla,
  alcance       = Catálogo de componentes,
  formulario    = Clase \texttt{lafrox} v2.2,
  fecha         = 17 de septiembre de 2026,
  version       = 2.2,
  nomenclatura  = SOBRE2\_LAFROX\_OFERTA\_TECNICA\_20261005.ZIP,
}

\begin{document}

\portadalafrox

% ---------------------------------------------------------------------------
\indicelafrox

% ===========================================================================
\chapter{Jerarquía, numeración y tabulaciones}
\label{cap:tipografia}

Todo el documento va en Inter a 11 puntos, el mínimo del Art. 40.4. Inter es
el sustituto abierto de la familia SF de Apple: mismo esqueleto
geométrico-humanista y altura de x alta. La jerarquía se construye con tamaño,
peso y espacio en blanco, sin filetes ni cajas, de modo que un subtítulo de
tercer nivel sigue distinguiéndose sin ocupar el espacio de un titular.

\nota{Las notas de orientación como esta se apagan de golpe compilando con la
opción de clase \texttt{[final]} o \texttt{[sinnotas]}: no hay que borrarlas
una por una antes de entregar. Aquí van también los pendientes ---por ejemplo,
que falta la volumetría de septiembre--- con un solo interruptor para todo el
andamiaje de redacción.}

\section{Niveles de título y su numeración}

Los capítulos numeran en árabe y arrastran la numeración de secciones, tablas
y figuras. La numeración llega hasta el tercer nivel y el índice, hasta el
segundo.

\subsection{Tercer nivel, para desagregar un requisito}

Este es un \texttt{subsection}. Por debajo hay un cuarto nivel numerado y un
quinto sin numerar, que corre en la misma línea del párrafo.

\subsubsection{Cuarto nivel, numerado}

Útil para separar los componentes de una decisión de arquitectura sin abrir
una sección nueva.

\paragraph{Quinto nivel} Corre dentro del párrafo y sirve para rotular un
argumento breve, como el contexto o la consecuencia de una decisión, sin
fragmentar la lectura.

\section{Tabulaciones}

\subsection{Listas con marca}

Tres niveles, cada uno con su propia marca y sangría:

\begin{itemize}
  \item La operación es dispersa y de terreno: preventa, reparto, preparación
        y recepción ocurren en la calle y en el local del cliente.
        \begin{itemize}
          \item 62 preventistas y cerca de 200 conductores.
          \item 14.200 puntos de entrega, con rutas rurales sin cobertura.
            \begin{itemize}
              \item Almacenes sin conexión a internet.
              \item Firma de recepción en papel.
            \end{itemize}
        \end{itemize}
  \item El perfil de carga no es plano: preparación nocturna, despacho
        concentrado y un septiembre que casi duplica el volumen.
\end{itemize}

\subsection{Listas numeradas}

Cuando el orden importa —etapas, pasos de un procedimiento, jerarquía de
precedencia— la numeración cambia de estilo por nivel:

\begin{enumerate}
  \item Bases Administrativas.
  \item Bases Técnicas Transversales.
    \begin{enumerate}
      \item Requisitos obligatorios.
      \item Requisitos deseables.
        \begin{enumerate}
          \item Con valor por defecto del transversal.
          \item Con valor del Capítulo 15 del caso.
        \end{enumerate}
    \end{enumerate}
  \item Caso 02 --- Logística. Puede endurecer un requisito transversal, nunca
        rebajarlo.
\end{enumerate}

\subsection{Listas de definición}

\begin{description}
  \item[OTIF] Entregas completas y a tiempo. Línea base 82,4\,\%.
  \item[Costo de servir] Costo total de atender un punto de entrega, incluidos
        reintentos y devoluciones.
  \item[Trazabilidad sanitaria] Capacidad de reconstruir la cadena de frío de
        un lote desde la recepción hasta la entrega.
\end{description}

\section{Citas textuales}

Para reproducir un pasaje del caso, de las Bases o de una entrevista, con la
fuente al pie:

\begin{cita}{Caso 02 --- Logística, Cap. 15}
El perfil de carga no es plano: septiembre casi duplica el volumen durante tres
semanas, y un dimensionamiento basado en el promedio está equivocado.
\end{cita}

El texto vuelve al cuerpo normal después de la cita, sin sangría heredada ni
cambio de interlínea.

% ===========================================================================
\chapter{Tablas y referencias a planillas}
\label{cap:tablas}

Las tablas van a 9,5 puntos, sobre el mínimo de 9 puntos que exige el
Art. 40.4, con cifras de ancho fijo para que las columnas de códigos y de
cantidades se alineen verticalmente. No hay relleno de cabecera, filas cebra
ni líneas verticales: tres filetes horizontales y espacio. La tabla rompe de
página sola y repite la cabecera.

Las columnas de ancho fijo se declaran en centímetros con \texttt{L\{\}},
\texttt{C\{\}} y \texttt{R\{\}}, y al menos una columna debe ser \texttt{Y}:
es la elástica, la que absorbe el resto del ancho para que la tabla cierre
exactamente en el margen.

\begin{tablalafrox}{Reparto de alcance entre etapas}{tab:etapas}%
  {L{2.1cm} Y C{2.4cm}}%
  {\cab{Etapa} & \cab{Alcance} & \cab{Producción}}
  Etapa 1 & Preventa, reparto, trazabilidad sanitaria y conteo cíclico
          & Mes 16 \\
  Etapa 2 & Costo de servir, planificación de rutas y analítica de OTIF
          & Mes 21 \\
  Operación & 36 meses de servicio gestionado, meses 21 a 56
          & --- \\
\end{tablalafrox}

Para los catálogos largos --- RF, RNF, Formulario T-12 --- la mecánica es la
misma; sólo cambia el número de filas:

\begin{tablalafrox}{Extracto del catálogo de requerimientos}{tab:catalogo}%
  {L{2.2cm} Y C{2.0cm} L{2.9cm}}%
  {\cab{Código} & \cab{Requerimiento} & \cab{Origen} & \cab{Componente}}
  RF-041 & Conteo cíclico con exactitud verificable por ubicación
         & Cap. 15.2 & WMS \\
  RF-042 & Registro de excursión térmica por lote y por tramo
         & RN-05 & IoT + WMS \\
  RNF-009 & Operación offline-first en rutas sin cobertura
         & R-12 & App de reparto \\
  RNF-010 & Ventana de despacho de 05:30 a 07:00 sin degradación
         & Cap. 15.4 & Núcleo transaccional \\
\end{tablalafrox}

\section{Referencias a planillas externas}

Cada subdocumento guarda las tablas grandes como planillas, una hoja por
tabla, dentro de su carpeta \texttt{entrega\_N/tablas/}. En el documento va la
referencia, no la tabla volcada a mano:

\tablaref{7}{Volumetría de despacho por ventana horaria}%
  {04\_arquitectura/entrega\_2/tablas/arquitectura\_fisica\_Dimensionamiento.xlsx!Despacho}

\tablaref{12}{Matriz de cumplimiento técnico RT-01 a RT-13}%
  {Productos/Formulario\_T-12\_Matriz\_Cumplimiento\_TFEP-01-2026.xlsx!T-12}

% ===========================================================================
\chapter{Páginas de control}
\label{cap:control}

\begin{controlversiones}
  \versionfila{0.1}{07-09-2026}{Patricio Henríquez}{Línea base de la Entrega 1}
  \versionfila{0.9}{29-09-2026}{Bastián Trejo}{Arquitectura lógica v6-2 incorporada}
  \versionfila{1.0}{05-10-2026}{Alex Aravena}{Versión presentada, Informe 2}
\end{controlversiones}

\begin{hojaresumen}
  \resumenfila{Jerarquía, numeración y tabulaciones}{cap:tipografia}
  \resumenfila{Tablas y referencias a planillas}{cap:tablas}
  \resumenfila{Páginas de control}{cap:control}
\end{hojaresumen}

\section{Verificación de cumplimiento formal}

\begin{tablalafrox}{Requisitos de forma del Art. 40 y su implementación}{tab:art40}%
  {C{1.3cm} Y L{4.7cm}}%
  {\cab{Art.} & \cab{Exigencia} & \cab{Dónde se resuelve}}
  40.1 & Foliación correlativa, sin saltos, extremo inferior derecho
       & \texttt{fancyhdr} en todos los estilos, portada incluida \\
  40.2 & Media firma en el extremo inferior derecho de cada página
       & Opción \texttt{[firma]} \\
  40.3 & Hoja resumen con folio inicial por sección
       & \texttt{hojaresumen} + \texttt{\textbackslash pageref} \\
  40.4 & Carta u oficio vertical
       & Opciones \texttt{[carta]} y \texttt{[oficio]} \\
  40.4 & Cuerpo no inferior a 11 pt; tablas y figuras, a 9 pt
       & \texttt{11pt} base; tablas a 9,5 pt \\
  40.4 & Índice detallado con folios
       & \texttt{\textbackslash indicelafrox} \\
  40.4 & Referencias en norma APA 7.ª ed.
       & Opción \texttt{[apa]} (\texttt{biblatex-apa}) \\
  40.2 & Firma \emph{completa} en la carátula
       & Bloque de firma de la portada, sobre papel blanco \\
  49.1 & Bloque de identificación que debe leerse en la carátula
       & Claves \texttt{objeto}, \texttt{sobre}, \texttt{correo} y
         \texttt{telefono} de \texttt{\textbackslash datoslafrox} \\
  50.3 & Nomenclatura del archivo
       & Clave \texttt{nomenclatura}, al pie de la carátula \\
\end{tablalafrox}

\nota{La opción \texttt{[borrador]} con la que se compiló este muestrario
imprime la marca de agua diagonal desde el folio 2. La carátula queda limpia a
propósito: la diagonal cruzaba el bloque de identificación del Art. 49.1, y el
estado del documento ya se declara dos veces en ella. Quita la opción --- o usa
\texttt{[final]} --- para la versión que se presenta.}

% ---------------------------------------------------------------------------
%  Cierre institucional. Comparte la diagonal de la carátula para que las dos
%  tapas se lean como un par.
\contraportadalafrox{%
  \begin{tablalafrox}{Equipo de trabajo del PROPONENTE}{tab:equipo}%
    {Y L{6.2cm}}%
    {\cab{Rol} & \cab{Nombre}}
    Jefe de Proyecto                        & Alex Aravena       \\
    Arquitecto de Solución                  & Bastián Trejo      \\
    Encargado de Seguridad de la Información & Álvaro Catalán    \\
    Líder de Datos                          & Leandro Chamorro   \\
    Líder de Desarrollo                     & Tomás Pérez        \\
    Líder de Calidad                        & Maximiliano Miño   \\
    Líder de Operación / SRE                & Guillermo Castillo \\
    Líder de Implantación y Gestión del Cambio & Patricio Henríquez \\
  \end{tablalafrox}}

\end{document}
